\begin{abstract}
We introduce \textbf{VEGA}, an approach for training navigation Vision-Language-Action (VLA) models from unlabeled egocentric navigation videos.
Internet-scale egocentric videos provide a scalable source of navigation-relevant visual observations, capturing cluttered scenes, close-range obstacles, and natural human motion through real-world spaces.
However, these videos are not directly usable for policy learning because they do not provide obstacle-aware trajectories conditioned on explicit navigation goals in the robot's coordinate frame.
VEGA addresses this gap by reconstructing local scene geometry from monocular video, sampling navigation goals (represented as text, image, or spatial waypoints) and generating obstacle-aware trajectories using the constructed geometry.
The resulting trajectory distribution is then used to train a flow-matching VLA navigation policy.
By using geometry exclusively during training, VEGA distills obstacle-aware planning directly into a vision-based policy. 
Furthermore, we introduce \textbf{VEGA-Bench}, a benchmark containing 250k scenes and approximately 8 million navigation goals paired with scene geometry, designed to evaluate goal progress, collision avoidance, and obstacle clearance of VLAs. 
In real-world trials, VEGA improves success by at least 65 percentage points, reduces collisions by at least 61.5\%, and improves obstacle clearance by at least 51.4\% relative to the evaluated baselines.
On VEGA-Bench, VEGA achieves competitive goal progress while reducing collisions by 18.2\% and improving obstacle clearance by 22.1\% over the strongest baseline under MoGe-2-based evaluation.
We further validate VEGA in dynamic navigation scenarios with moving pedestrians and static obstacles, providing additional evidence of obstacle-aware navigation under dynamic conditions.
VEGA achieves these results with supervision from just 250 hours of egocentric video, approximately 2.6\% of the strongest navigation baseline's training-data duration, without collecting robot demonstrations.
This demonstrates the scalability and effectiveness of video-derived geometric supervision for training obstacle-aware navigation VLAs.
% The code and benchmark will be released at the time of publication.
\end{abstract}

\keywords{Vision-Language-Action Models, Learning from Video, Geometric Supervision}

\vspace{-2ex}
% Two or three meaningful keywords should be added here

% Our approach is motivated by the fact that internet-scale egocentric videos provide a scalable source of navigation-relevant visual observations across diverse real-world environments including cluttered scenes and close-range obstacles. 
